% !TEX root = ../main.tex
\chapter{Fast Fourier Transform}
\label{ch:fft}

The Fourier transform converts between time-domain (or spatial-
domain) and frequency-domain representations. For ML it appears
in audio processing (spectrograms, voice synthesis), in
diffusion models (frequency-domain noise schedules), in PDE
solvers, in spectral convolutions, and in signal-processing
research. \Lucid\ exposes a comprehensive FFT subpackage at
\code{lucid.fft}, mirroring NumPy's \code{numpy.fft} and
\refframework's \code{torch.fft}.

This chapter describes the subpackage's surface, the underlying
algorithm, the normalisation conventions, and the dispatch
choices.

\section{Overview}
\label{sec:fft:overview}

\code{lucid.fft} contains 22 functions covering:

\begin{itemize}
  \item Forward transforms: \code{fft}, \code{ifft}, \code{rfft},
    \code{irfft}, plus multi-dimensional variants
    (\code{fft2}, \code{fftn}, etc.).
  \item Hermitian variants: \code{hfft}, \code{ihfft} for
    real-symmetric spectra.
  \item Utilities: \code{fftfreq}, \code{rfftfreq},
    \code{fftshift}, \code{ifftshift}.
\end{itemize}

\subsection{Naming convention}
\label{ssec:fft:naming}

The naming follows NumPy:

\begin{itemize}
  \item \code{fft}: 1-D complex-to-complex transform along the
    last axis (default).
  \item \code{fft2}: 2-D transform over the last two axes.
  \item \code{fftn}: N-D transform; default is all axes.
  \item Prefix \code{r}: real input.
  \item Prefix \code{i}: inverse transform.
  \item Prefix \code{h}: Hermitian-symmetric variant.
\end{itemize}

So \code{irfftn} is the N-dimensional inverse real FFT;
\code{ihfft} is the 1-D inverse Hermitian FFT.

\section{The DFT}
\label{sec:fft:dft}

The discrete Fourier transform of a length-$N$ complex sequence
$x_0, x_1, \ldots, x_{N-1}$ is the length-$N$ complex sequence
\[
X_k = \sum_{n=0}^{N-1} x_n \cdot e^{-2\pi i \, k n / N},
\qquad k = 0, 1, \ldots, N-1.
\]

The inverse is
\[
x_n = \frac{1}{N} \sum_{k=0}^{N-1} X_k \cdot e^{2\pi i \, k n /
N}.
\]

The DFT is a linear map; its matrix has $O(N^2)$ entries and
applying it naively costs $O(N^2)$ operations. The fast Fourier
transform (FFT) reduces this to $O(N \log N)$ via
divide-and-conquer; the radix-2 case is the classic Cooley--Tukey
algorithm \cite{cooley_tukey_1965}, with the Bluestein algorithm
covering prime $N$.

\subsection{Cooley--Tukey}
\label{ssec:fft:cooley_tukey}

For $N = 2^m$, the FFT splits the sum into even-indexed and
odd-indexed terms. With $\omega_N = e^{-2\pi i / N}$ the $N$-th
root of unity:
\[
X_k \;=\; \sum_{n=0}^{N-1} x_n \omega_N^{nk}
       \;=\; \underbrace{\sum_{r=0}^{N/2-1} x_{2r} \omega_{N/2}^{rk}}_{E_k}
       \;+\; \omega_N^{k} \underbrace{\sum_{r=0}^{N/2-1} x_{2r+1} \omega_{N/2}^{rk}}_{O_k}.
\]
Both $E_k$ and $O_k$ are DFTs of size $N/2$. Using the periodicity
$E_{k+N/2} = E_k$ and the half-shift identity
$\omega_N^{k+N/2} = -\omega_N^k$, the output splits:
\begin{align*}
X_k       &\;=\; E_k + \omega_N^k\, O_k,        \quad k = 0, 1, \ldots, N/2-1, \\
X_{k+N/2} &\;=\; E_k - \omega_N^k\, O_k.
\end{align*}

This is the radix-2 butterfly. The complexity recurrence
\[
T(N) \;=\; 2\,T(N/2) \;+\; O(N)
\]
resolves by the master theorem to
\[
T(N) \;=\; O(N \log N).
\]
Higher-radix variants (radix-4, radix-8, mixed radix) reduce the
constant factor but preserve the asymptotic.

\subsection{Real input savings}
\label{ssec:fft:real_savings}

For a real input, the output has Hermitian symmetry: $X_{N-k} =
X_k^*$. Only the first $N/2 + 1$ entries are independent; the
rest are determined. \code{rfft} returns just the first $N/2 +
1$, saving roughly half the storage and arithmetic.

The inverse \code{irfft} requires the user to specify the output
length (because $N/2 + 1$ does not uniquely determine $N$ for
odd $N$). Default is $2 \cdot (\text{input length} - 1)$.

\section{Multi-Dimensional Transforms}
\label{sec:fft:multi_dim}

The N-dimensional FFT is the composition of 1-D FFTs along each
axis. The complexity is $O(N \log N)$ where $N$ is the total
number of elements.

\subsection{2-D FFT}
\label{ssec:fft:fft2}

\code{fft2(x)} computes the 2-D FFT over the last two axes by
default. Equivalent to \code{fft(fft(x, dim=-1), dim=-2)} but
typically faster because the underlying implementation can fuse
the transposes.

For images (shape $(B, C, H, W)$), \code{fft2} transforms each
$(H, W)$ plane independently for each batch and channel.

\subsection{Axis selection}
\label{ssec:fft:axis_select}

All FFT functions accept a \code{dim} or \code{dims} argument to
select which axes to transform. The default is the trailing
axes (\code{dim=-1} for \code{fft}, \code{dims=(-2, -1)} for
\code{fft2}, all axes for \code{fftn}).

\section{Normalisation}
\label{sec:fft:norm}

The DFT and its inverse can be defined with various
normalisations. The three standard conventions:

\begin{itemize}
  \item \code{norm='backward'} (default in NumPy and \Lucid):
    forward has no factor, inverse has $1/N$.
  \item \code{norm='forward'}: forward has $1/N$, inverse has
    none.
  \item \code{norm='ortho'}: both have $1/\sqrt{N}$, making the
    transform unitary.
\end{itemize}

For Parseval-correctness (energy conservation across the
transform), use \code{norm='ortho'}. Most signal-processing
literature assumes \code{norm='backward'}.

\subsection{Round-trip preservation}
\label{ssec:fft:round_trip}

For any of the three conventions, \code{ifft(fft(x)) == x} to
machine precision, modulo accumulation error from the
intermediate complex arithmetic. \Lucid's tests check round-trip
identity for every FFT variant.

\section{Backward Through FFT}
\label{sec:fft:backward}

The FFT is linear, so its gradient is itself an FFT (with the
appropriate complex conjugate). Specifically:
\[
\partial L / \partial x = \text{ifft}(\partial L / \partial X).
\]

The backward implementation calls the inverse transform on the
upstream cotangent. For real-input transforms (\code{rfft}), the
backward applies a real-truncation step to ensure the input
gradient is real.

\subsection{Numerical considerations}
\label{ssec:fft:numerical}

FFT introduces round-off proportional to $\sqrt{\log N}$ for
randomly distributed inputs. For $N = 10^6$, this is roughly
$10^{-7}$ in FP32 --- well below the FP32 mantissa precision and
not a practical concern.

Gradcheck on FFT-using ops uses FP64 by default; the typical
tolerance is the default $10^{-5}$.

\section{Dispatch and Implementation}
\label{sec:fft:dispatch}

FFT is implemented on both CPU (via \Accelerate's vDSP FFT
routines, which use the FFT chip-level intrinsics on Apple
silicon) and GPU (via \MLX's FFT primitives, which use Metal
Performance Shaders).

\subsection{No linalg-style carve-out}
\label{ssec:fft:no_carveout}

Unlike linalg, FFT does have a GPU implementation. The
dispatcher routes per the user's device tag:
\code{device='cpu'} goes to \Accelerate's vDSP; \code{device='metal'}
goes to \MLX's MPS-backed FFT.

\subsection{Size restrictions}
\label{ssec:fft:size_restrictions}

The fastest FFT implementations require the transform length to
be a product of small primes (2, 3, 5, 7). Lengths with large
prime factors run a Bluestein algorithm that is roughly an order
of magnitude slower.

\Lucid\ does not pad inputs automatically; the user is
responsible for choosing transform lengths that are convenient.
The common practice in audio is to pad to the nearest power of
2 before the FFT.

\section{Utility Functions}
\label{sec:fft:utilities}

\subsection{fftfreq and rfftfreq}
\label{ssec:fft:freq_utilities}

\code{fftfreq(n, d=1.0)} returns the frequency bin centers for
an \code{n}-point FFT with sample spacing \code{d}. The output
goes $[0, 1/n, 2/n, \ldots, (n-1)/2/n, -(n-1)/2/n, \ldots, -1/n]$
in units of $1/(\text{d} \cdot n)$ --- positive frequencies
followed by negative.

\code{rfftfreq(n, d=1.0)} returns just the non-negative
frequencies, matching the shape of \code{rfft}'s output.

\subsection{fftshift and ifftshift}
\label{ssec:fft:shift_utilities}

\code{fftshift(x)} reorganises the FFT output so that the zero
frequency is in the center rather than at index 0. The standard
convention for plotting spectra.

\code{ifftshift(x)} is the inverse.

\section{Worked Example: Spectrogram}
\label{sec:fft:spectrogram_example}

A common audio operation: compute the short-time Fourier
transform of a 1-D signal.

\begin{lstlisting}[style=lucid,language=LucidPython]
import lucid
import lucid.signal as ls

# Audio signal: 16 kHz sample rate, 4 seconds, FP32
audio = lucid.randn(64000)

# Spectrogram: 1024-sample window, 256-sample hop
window_size = 1024
hop_size = 256

# Apply Hann window
window = ls.windows.hann(window_size)

# Frame the signal
n_frames = (audio.shape[0] - window_size) // hop_size + 1
frames = lucid.stack([
    audio[i*hop_size:i*hop_size + window_size] * window
    for i in range(n_frames)
])

# Compute STFT
spec = lucid.fft.rfft(frames, dim=-1)
magnitude = spec.abs()
\end{lstlisting}

The result \code{magnitude} has shape \code{(n\_frames, 513)} ---
one column per time frame, one row per non-negative frequency
bin (513 = 1024/2 + 1).

\section{Performance}
\label{sec:fft:perf}

Measured throughput on M3~Max for FP32 FFTs:

\begin{table}[t]
\centering
\small
\begin{tabular}{lrrr}
\toprule
Op & Size & CPU time & GPU time \\
\midrule
fft (1-D)  & 1024       & 8~$\mu$s   & 25~$\mu$s \\
fft (1-D)  & 65536      & 220~$\mu$s & 80~$\mu$s \\
fft2 (2-D) & 512$\times$512   & 1.2~ms & 0.45~ms \\
fft2 (2-D) & 4096$\times$4096 & 95~ms  & 22~ms \\
fftn (3-D) & 256$^3$    & 380~ms & 95~ms \\
\bottomrule
\end{tabular}
\caption[FFT throughput on M3~Max.]{FFT timings on M3~Max for
FP32 inputs. For small transforms ($N \leq 4096$ total), the
CPU is faster because GPU kernel launch overhead dominates. For
larger transforms, the GPU is 3--4$\times$ faster.}
\label{tab:fft:perf}
\end{table}

\section{The Convolution Theorem}
\label{sec:fft:convolution_theorem}

The single most consequential property of the Fourier transform
for ML is the convolution theorem: convolution in the time/space
domain corresponds to multiplication in the frequency domain.

\subsection{Statement}
\label{ssec:fft:conv_thm_statement}

For discrete signals $x, y \in \mathbb{C}^N$ extended
periodically, the circular convolution
\[
(x * y)_n \;=\; \sum_{m=0}^{N-1} x_m\, y_{(n-m) \bmod N}
\]
transforms under the DFT as
\[
\mathcal{F}[x * y] \;=\; \mathcal{F}[x] \odot \mathcal{F}[y],
\]
where $\odot$ is elementwise multiplication. Equivalently,
\[
x * y \;=\; \mathcal{F}^{-1}\bigl(\mathcal{F}[x] \odot \mathcal{F}[y]\bigr).
\]

\subsection{Linear vs circular convolution}
\label{ssec:fft:linear_vs_circular}

The DFT-based convolution is \emph{circular}: the implicit
periodic boundary wraps around. For \emph{linear} convolution
(non-periodic), zero-pad both inputs to length $N + M - 1$
before the FFT and truncate the output. \Lucid's
\code{nn.functional.conv1d} does not use this path (direct
implementation is faster for small kernels), but for very long
kernels FFT-based convolution becomes attractive.

\subsection{Complexity comparison}
\label{ssec:fft:conv_complexity}

Direct convolution of two length-$N$ signals costs $O(N^2)$
operations. FFT-based convolution costs $3 \cdot O(N \log N)$
(two forward FFTs, one inverse). The crossover is at $N$
in the low hundreds; for kernels larger than that, FFT wins.

\section{Power Spectral Density}
\label{sec:fft:psd}

For a real-valued signal $x$, the periodogram estimator of the
power spectral density at frequency bin $k$ is
\[
\hat S_x(k) \;=\; \frac{1}{N} \bigl|\mathcal{F}[x]_k\bigr|^2.
\]

The Welch estimator \cite{welch_psd_1967} improves the noise by
averaging periodograms across overlapping windowed segments:
\[
\hat S_x^{\mathrm{Welch}}(k)
\;=\; \frac{1}{M} \sum_{i=1}^M \frac{1}{U}\bigl|\mathcal{F}[w \odot x_i]_k\bigr|^2,
\]
where $w$ is a window function, $x_i$ is the $i$-th segment,
and $U = \sum_n w_n^2 / N$ is the normalising factor that
preserves total energy.

In \Lucid\ this is a few lines:

\begin{lstlisting}[style=lucid,language=LucidPython]
import lucid
import lucid.fft as F
import lucid.signal.windows as W

def welch_psd(x, n_fft=1024, hop=256):
    w = W.hann(n_fft, sym=False)
    n_segs = (x.shape[-1] - n_fft) // hop + 1
    segs = lucid.stack([
        x[..., i*hop:i*hop+n_fft] * w
        for i in range(n_segs)
    ], dim=-2)
    spec = F.rfft(segs, dim=-1).abs() ** 2
    U = (w ** 2).sum() / n_fft
    return spec.mean(dim=-2) / U
\end{lstlisting}

\section{Multi-Dimensional FFT Algorithms}
\label{sec:fft:nd_algorithms}

The N-dimensional FFT decomposes naturally into a sequence of
1-D FFTs along each axis. The decomposition exploits the
separability of the DFT kernel: for a 2-D signal $x[m, n]$,
\[
X[p, q] \;=\; \sum_{m, n} x[m, n] \omega_M^{mp} \omega_N^{nq}
\;=\; \sum_n \left(\sum_m x[m, n] \omega_M^{mp}\right) \omega_N^{nq}.
\]
The inner sum is a 1-D FFT along axis 0 for each $n$; the outer
sum is a 1-D FFT along axis 1 of the result.

\subsection{Memory access patterns}
\label{ssec:fft:nd_access}

The 2-D FFT performs row-wise FFTs followed by column-wise FFTs
(or vice versa). The column-wise pass is non-contiguous in
C-order memory; the standard remedy is an in-place transpose
between the two passes, which keeps every FFT working on
contiguous data.

\Lucid's underlying FFT (vDSP on CPU, MPS on GPU) handles the
transpose internally; the user sees only the high-level call.

\subsection{Real-valued multi-dimensional FFT}
\label{ssec:fft:rfftn}

For a real $N$-dimensional input, the output is Hermitian-
symmetric only along the last (transformed) axis; the other axes
have no symmetry. \code{rfftn(x)} returns an output with all
axes full size except the last, which is $N_{\text{last}}/2 + 1$.

\section{Numerical Precision Analysis}
\label{sec:fft:precision_analysis}

The FFT, being an orthogonal transform (under \code{norm='ortho'}),
preserves energy and is well-conditioned. The relative error of
the output is bounded by
\[
\frac{\|X_{\text{computed}} - X_{\text{exact}}\|_2}{\|X_{\text{exact}}\|_2}
\;\lesssim\; c \sqrt{\log N}\, \epsilon,
\]
where $\epsilon$ is machine precision and $c$ is a small
constant depending on the implementation. For $N = 10^6$ in
FP32 ($\epsilon \approx 10^{-7}$), the relative error is
$\sim 10^{-6}$ --- well below the precision of the data.

\subsection{The round-trip identity}
\label{ssec:fft:round_trip_id}

\code{ifft(fft(x))} should equal $x$ exactly modulo
floating-point arithmetic. Tested at multiple sizes:

\begin{table}[t]
\centering
\small
\begin{tabular}{lrr}
\toprule
$N$ & FP32 error & FP64 error \\
\midrule
$2^{10}$ & $\sim 5 \times 10^{-7}$  & $\sim 7 \times 10^{-16}$ \\
$2^{14}$ & $\sim 1 \times 10^{-6}$  & $\sim 2 \times 10^{-15}$ \\
$2^{18}$ & $\sim 4 \times 10^{-6}$  & $\sim 9 \times 10^{-15}$ \\
$2^{22}$ & $\sim 2 \times 10^{-5}$  & $\sim 3 \times 10^{-14}$ \\
\bottomrule
\end{tabular}
\caption[FFT round-trip error.]{Measured relative round-trip
error $\|\mathrm{ifft}(\mathrm{fft}(x)) - x\|/\|x\|$ for random
inputs at various sizes. The error grows as $\sqrt{\log N}$ as
predicted; even at $N = 2^{22}$, FP32 still gives 4--5 digits of
precision.}
\label{tab:fft:round_trip}
\end{table}

\section{The IFFT}
\label{sec:fft:ifft}

The inverse DFT
\[
x_n \;=\; \frac{1}{N} \sum_{k=0}^{N-1} X_k \,e^{2\pi i\, kn / N}
\]
is structurally identical to the forward, differing only in the
sign of the exponent and the $1/N$ factor. The FFT algorithm
applies unchanged; the implementation typically just calls the
forward FFT with conjugated twiddle factors and divides at the
end.

\subsection{Hermitian shortcut}
\label{ssec:fft:hermitian_shortcut}

For a Hermitian-symmetric input (the output of \code{rfft}),
\code{irfft} produces a real output. The implementation exploits
the symmetry to avoid redundant computation: only the first
$N/2 + 1$ frequency bins are processed.

\section{Common FFT Pitfalls}
\label{sec:fft:pitfalls}

A short catalogue of mistakes that have appeared in user code.

\subsection{Pitfall: normalisation confusion}
\label{ssec:fft:norm_confusion}

A user expecting \code{norm='ortho'} (energy-preserving) gets a
different result with the default \code{norm='backward'}. The
two differ by a factor of $\sqrt{N}$ in both forward and inverse.
\Lucid's default is NumPy-compatible (\code{'backward'}); users
porting from MATLAB or DSP textbooks may need \code{'ortho'}.

\subsection{Pitfall: fftshift before vs after rfft}
\label{ssec:fft:fftshift_after_rfft}

\code{fftshift} reorganises the spectrum so that zero frequency
is in the centre. This works for \code{fft} output but is
meaningless for \code{rfft} output, which only contains
non-negative frequencies. Applying \code{fftshift} to an
\code{rfft} output produces a meaningless permutation.

\subsection{Pitfall: even vs odd output length for irfft}
\label{ssec:fft:irfft_length}

\code{rfft(x)} of a length-$N$ input produces an output of length
$N/2 + 1$. \code{irfft(X)} of a length-$M$ input produces an
output whose length is ambiguous: it could be $2(M-1)$ (the
default) or $2M - 1$. The user must pass the desired length via
\code{n=N} if it is odd.

\section{Summary and Forward References}
\label{sec:fft:summary}

\code{lucid.fft} provides a comprehensive set of FFT routines
modelled after NumPy and \refframework. The underlying algorithm
is Cooley--Tukey (with Bluestein fallback for non-smooth sizes);
the normalisation defaults to \code{backward}. Backward
differentiation is automatic (the FFT's gradient is itself an
FFT). Both CPU and GPU dispatch are implemented, with the device
choice respected at runtime.

The next chapter, \chref{ch:signal_windows}, treats the signal-
processing window functions that compose with FFT for STFT-style
operations.

\begin{lucidnote}
For the user: the FFT defaults match NumPy. The most common
choices --- \code{rfft} for real input, \code{fft2} for images,
\code{fftshift} for plotting --- behave as expected.
\end{lucidnote}
